\selectlanguage{ngerman}%
\frenchspacing%
\begin{ledgroupsized}[r]{120mm}%
\footnotesize%
\pstart%
\noindent\textbf{Überlieferung:}%
\pend%
\end{ledgroupsized}%
\begin{ledgroupsized}[r]{114mm}%
\footnotesize%
\pstart%
\parindent -6mm%
\makebox[6mm][l]{\textit{L}}%
Aufzeichnung:
LH~XXXVII,\;5 Bl.\;171--172.
Ein Bogen, 8\textsuperscript{o},
Ränder zum Teil beschnitten;
Papiererhaltungsmaßnahmen.
Fragment eines \href{\Scan{LH_37_05_171-172wz}}{Wasserzeichens} am oberen Rand.
Eine ganzseitig beschriebene Seite auf Bl.\;\href{\Scan{LH_37_05_171r}}{171\,r\textsuperscript{o}}.
Bl.\;\href{\Scan{LH_37_05_171v}}{171\,v\textsuperscript{o}} und
Bl.\;\href{\Scan{LH_37_05_172r}}{172\,r\textsuperscript{o}} sind unbeschrieben.
Bl.\;\href{\Scan{LH_37_05_172v}}{172\,v\textsuperscript{o}} überliefert die Aufzeichnung N.\;\ref{60349-2}.
\pend%
\end{ledgroupsized}%
\vspace{5mm}%
\begin{ledgroup}%
\footnotesize%
\pstart%
\noindent%
\textbf{Datierungsgründe:}
Die Datierung der Aufzeichnung N.\;\ref{60349-1} beruht auf Eigenschaften des Textträgers sowie auf inhaltlichen Betrachtungen.
Siehe die ausführliche Begründung auf S.\;\pageref{60349-1_dat}.%
\pend%
\end{ledgroup}%
\selectlanguage{latin}%
\frenchspacing%
\vspace{8mm}%
\normalsize%
\pstart%
\noindent%
\bracks{%
\href{\Scan{LH_37_05_171r}}%
{171\,r\textsuperscript{o}}%
}
Mihi videtur
\edtext{%
non quies%
\indexS{quies causa unionis corporum}
esse causa unionis corporum%
\indexS{causa unionis corporum}%
\indexS{unio corporum}%
}{%
\lemma{quies \bracks{...} corporum}%
\Cfootnote{%
Anspielung auf die kartesische Erklärung der Körperkohäsion.
Siehe etwa \textsc{R.\;Descartes},
\textit{Principia philosophiae} pars~II, §\;54--55
(Amsterdam 1644, S.\;62\,f.;\cite{00035}
\textit{DO} VIII,\;1, S.\;70\,f.).\cite{00120}%
}}
sed motus.%
\indexS{motus causa unionis}
Motus inquam corporum ambientium,%
\indexS{motus corporum ambientium}%
\indexS{corpora ambientia}
quia duo
\edtext{%
corpora sine turbatione ejus%
\indexS{turbatio motus e disjuntione corporum}
disjungi%
\indexS{corpora disjuncta}%
\indexS{disjunctio corporum}%
}{%
\lemma{corpora}%
\Bfootnote{%
\textit{(1)}~disjungi %
\textit{(2)}~sine turbatione ejus disjungi%
~\textit{L}%
}}
non possunt,
et quo majus turbatio est,%
\indexS{turbatio motus e disjuntione corporum}
eo major est resistentia.%
\indexS{resistentia disjunctioni corporum}%
\indexS{resistentia e turbatione motus}
Ut fit
\edtext{%
cum pars fluidi%
\indexS{fluidum expulsum}%
\indexS{pars fluidi expulsa}
loco%
\indexS{locus fluidi expulsi}
expellitur%
\indexS{expulsio fluidi}%
\bracks{,}
alia ejusdem fluidi%
\indexS{fluidum succedens}
parte non succedente;%
\indexS{pars fluidi succedens}%
\indexS{successio fluidi}%
}{%
\lemma{cum}%
\Bfootnote{%
\textit{(1)}~corpus %
\textit{(a)}~li\angles{q} %
\textit{(b)}~fluidum %
\textit{(2)}~pars fluidi loco expellitur %
\textit{(a)}~nec %
\textit{(aa)}~ali\angles{ud} %
\textit{(bb)}~alia simile\bracks{\textit{!}} in eum succedere statim potest %
\textit{(b)}~alia ejusdem %
\bracks{...} non succedente;%
~\textit{L}%
}}
inde enim magnam mutationem sequi%
\indexS{mutatio sequens}
necesse est,
ut videmus
cum aquam antliis haurimus%
\indexS{aqua hausta antliis}%
\indexS{haustus aquae antliis}%
\indexS{antlia}
aut duo marmora polita%
\indexS{marmora polita}
a se invicem avellimus.%
\indexS{avulsio marmorum politorum}
Quando autem aer facile circulare potest,%
\indexS{aer circulans}%
\indexS{circulatio aeris}
ita ut
anteriore pulso%
\indexS{aer pulsus}
alius a tergo succedat,%
\indexS{aer succedens a tergo}
resistentia exigua est,%
\indexS{resistentia pulsioni aeris}
quia mutatio etiam exigua est,%
\indexS{mutatio sequens}
quae non in alio constabit,
quam quod corpus%
\indexS{corpus commovens aerem}
\edtext{%
motum aerem commovit,%
\indexS{aer commotus}
quantum satis erat ad circulationem,%
\indexS{circulatio aeris}
id est portionem ejus%
\indexS{portio aeris}
non multo mole majorem corpore moto.%
\indexS{corpus motum}%
}{%
\lemma{motum}%
\Bfootnote{%
\textit{(1)}~tantundem %
\textit{(2)}~tantu %
\textit{(3)}~aequalem suae aeris molem loco exp\angles{rimit}, %
\textit{(4)}~\angles{comm} %
\textit{(5)}~aerem commovit, %
\bracks{...} id est %
\textit{(a)}~\angles{partem} %
\textit{(b)}~portionem ejus %
\bracks{...} corpore moto.%
~\textit{L}%
}}
\pend%
\pstart%
Dum enim
\edtext{%
corpus \textit{A} transit in locum \textit{B},%
\indexS{corpus transiens in locum alium}%
}{%
\lemma{corpus \textit{A} \bracks{...} locum \textit{B}}%
\Cfootnote{%
Vgl. das Diagramm \bracks{\textit{Fig.\:1}}.%
}}
sufficit
\edtext{%
moveri aerem \textit{clBmf},%
\indexS{aer commotus}%
}{%
\lemma{moveri aerem}%
\Bfootnote{%
\textit{(1)}~\textit{fedc} %
\textit{(2)}~\angles{\textit{clBmf}}, %
\textit{(3)}~\textit{clBmf},%
~\textit{L}%
}}
id est aerem loci occupandi \textit{B} %
cum filo tenui \textit{cl.fm}%
\indexS{filum aeris tenue}
corpus \textit{A} radente;
nam corpore \textit{A} vel minimum progrediente
\edtext{%
versus~\textit{B},
nempe in locum \textit{medl}%
\indexS{locus occupandus}
deserto loco \textit{fghc},%
\indexS{locus desertus}
statim filum aeris \textit{emf}%
\indexS{filum aeris tenue}
ruit in \textit{mfg},
et filum \textit{dlc}
ruit in \textit{lch},%
\indexS{filum aeris tenue}
atque ita circulatio%
\indexS{circulatio aeris}
ad implendum locum desertum%
\indexS{locus desertus}%
\indexS{locus implendus}
necessaria,%
\indexS{circulatio necessaria ad locum implendum}
fit
quam minima aeris ambientis mutatione.%
\indexS{mutatio aeris ambientis}%
\indexS{aer ambiens}%
}{%
\lemma{versus~\textit{B},}%
\Bfootnote{%
\textit{(1)}~aer lateralis \textit{lc}, \textit{mf} versus \textit{ch}, \textit{fg} ruens,\indexS{aer ruens} in locum relictum\indexS{locus relictus} succedit, quem alius aer ex \textit{B} veniens %
\textit{(a)}~continue sequi %
\textit{(b)}~contr %
\textit{(c)}~et his filamentis\indexS{filamentum} velut adductis continuis supplementis\indexS{supplementum continuum} sequitur. Nam dum %
\textit{(aa)}~\textit{mf} %
\textit{(bb)}~\textit{lc} et \textit{mf} tendunt versus \textit{gh} %
\textit{(aaa)}~\angles{\textit{em}} %
\textit{(bbb)}~interea %
\textit{(aaaa)}~\angles{et} %
\textit{(bbbb)}~\textit{dl} et \textit{em} succedunt in locum %
\textit{(aaaaa)}~\textit{mf} %
\textit{(bbbbb)}~\textit{lc} et \textit{mf} eodemque modo corpus \textit{A} radunt. %
\textit{(2)}~nempe in %
\bracks{...} in \textit{mfg}, et %
\textit{(a)}~\textit{dlc} in \textit{lc} %
\textit{(b)}~filum \textit{dlc} %
\bracks{...} quam minima %
\textit{(aa)}~fluidi %
\textit{(bb)}~aeris ambientis mutatione.%
~\textit{L}%
}}
Unde
\edtext{%
si piscis moveatur in media aqua,%
\indexS{piscis movens sine turbatione aquae}
videmus tranquillitatem aquae%
\indexS{tranquillitas aquae}%
\indexS{aqua tranquilla}
non turbari notabiliter.%
\indexS{turbatio tranquillitatis aquae}%
}{%
\lemma{si}%
\Bfootnote{%
\textit{(1)}~\angles{cor}pus %
\textit{(2)}~piscis moveatur in media aqua, %
\textit{(a)}~ut piscis, %
\textit{(b)}~n\angles{ec} %
\textit{(c)}~\angles{et fundo} %
\textit{(d)}~et a fundo et superficie sic satis distet, tranquillitas aquae notabi %
\textit{(e)}~videmus tranquillitatem %
\textit{(aa)}~ejus %
\textit{(bb)}~no %
\textit{(cc)}~aquae non turbari notabiliter.%
~\textit{L}%
}}
Et quo
\edtext{%
magis medium%
}{%
\lemma{magis}%
\Bfootnote{%
\textit{(1)}~corpus %
\textit{(2)}~medium%
~\textit{L}%
}}
liquidum et homogeneum est,%
\indexS{medium liquidum}%
\indexS{medium homogeneum}
eo minus motui corporum resistit,%
\indexS{medium resistens motui}%
\indexS{resistentia medii}%
\indexS{motus corporum}
praesertim si et
\edtext{%
corpora,
aptissima ad medium dividendum ratione%
\indexS{ratio motus apta ad medium dividendum}%
}{%
\lemma{corpora}%
\Bfootnote{%
\bars{in r \textit{erg. u. gestr.}}%
, aptissima ad medium dividendum ratione%
~\textit{L}%
}}
moveantur.%
\indexS{corpora mota ratione apta ad medium dividendum}
\pend%
\begin{edfigure}[2.50em][1.50em]%
\includegraphics%
{gesamttex/edit_VIII-4/images/LH_37_05_171-172_d_171r.pdf}%
\edcaption[0.5em]{%
\bracks{\textit{Fig.\:1}}%
\edtext{%
}{%
\lemma{\bracks{\textit{Fig.\:1}}\,}%
\Cfootnote[nonum]{%
Zusätzlich zum Diagramm sind zwei gestrichene, nicht wiedergegebenen Entwürfe in \textit{L} (Bl.\;\href{\Scan{LH_37_05_171r}}{171\,r\textsuperscript{o}}) überliefert.%
}%
}%
}%
\end{edfigure}%
\vspace{4.0em}%
\pstart%
\noindent%
\centering%
*     *     *
\pend%
\newpage%
\selectlanguage{german}%
\frenchspacing%
\footnotesize%
\pstart%
\noindent%
\textls{Ausführliche Datierungsgründe:}%
\pagelabel{60349-1_dat}
Die undatierte und titellose Aufzeichnung N.\;\ref{60349-1} ist dem grundlegenden Phänomen der Körperkohäsion gewidmet.
Ausgangspunkt ist René Descartes'\indexN{\textls{Descartes} (Cartesius, des Cartes), René 1596\textendash1650} Erklärung, die relative Ruhe der Körperteile begründe deren Zusammenhalt und somit die Kohäsion der Körper (siehe den Quellennachweis im Erläuterungsapparat).
Über diese Erklärung hat sich Leibniz seit den frühen 1670er Jahren kritisch geäußert, etwa in der Aufzeichnung \textit{De consistentia corporum}, die auf die Monate zwischen Herbst und Winter 1672/1673 datierbar ist (\textit{LSB} VI,\;3 N.\;4, bes. S.\;95.6--18);\cite{01299}
seine Abneigung ihr gegenüber lässt sich auf die bereits in der \textit{Theoria motus abstracti} %
aufgestellte Annahme zurückführen, dass Körperkohäsion letztlich auf Bewegung beruht, nicht auf Ruhe (Mainz 1671; \textit{LSB} VI,\;2 N.\;41, S.\;269.28\,f.).\cite{00259}
In der vorliegenden Aufzeichnung N.\;\ref{60349-1} stellt Leibniz der Auffassung Descartes'\indexN{\textls{Descartes} (Cartesius, des Cartes), René 1596\textendash1650} die These entgegen, dass die Kohäsion der Teile eines Körpers von der Bewegung eines umgebenden Fluidums herrühre.
Verdeutlicht wird diese These am Beispiel eines Wurfels, der sich in der Luft fortbewegt.
\pend%
\pstart%
Der Träger der Aufzeichnung N.\;\ref{60349-1} weist am oberen Rand das Fragment eines \href{\Scan{LH_37_05_171-172wz}}{Wasserzeichens} auf, das von der Papiermühle bei Lauenstein\indexO{Lauenstein (Niedersachsen)} (Niedersachsen)\indexO{Niedersachsen} %
verwendet wurde und im Leibniz-Nachlass nach heutigem Kenntnisstand für die Jahre 1680 bis 1683 bzw. 1687 belegt ist. %
Die Aufzeichnung N.\;\ref{60349-1} ist demzufolge aller Wahrscheinlichkeit nach innerhalb dieses Zeitraums entstanden, %
der hier insgesamt als Datierung vorgeschlagen wird.
Diese Entscheidung lässt sich folgendermaßen begründen:
\pend%
\pstart%
Leibnizens Ablehnung von Descartes'\indexN{\textls{Descartes} (Cartesius, des Cartes), René 1596\textendash1650} Erklärung der Körperkohäsion, gekoppelt mit der These, dass die "erste Ursache" in der \textit{pressio ambientium} liege, äußert Leibniz auch in einer Auzeichnung, die er selbst \textit{Cohaesio} betitelt hat (\textit{LSB} VIII,\;3 N.\;16, S.\;262.7).\cite{02615}
Diese lässt sich aufgrund ihrer inhaltlichen Verwandschaft mit den Entwürfen \textit{Explicatio mechanica elastri} (\textit{LSB} VIII,\;3 N.\;14\textsubscript{3})\cite{01437} und \textit{Demonstratio quod extensiones elasticorum sint viribus tendentibus proportionales} (\textit{LSB} VIII,\;3 N.\;14\textsubscript{7})\cite{01438} auf die Zeitspanne von Ende Januar 1683 und der ersten Hälfte 1684 datieren (vgl. die Begründung in \textit{LSB} VIII,\;3 N.\;16, S.\;261).\cite{02615} 
Mit diesen Texten -- vornehmlich mit \textit{Cohaesio} -- ist auch die vorliegende Aufzeichnung N.\;\ref{60349-1} inhaltlich verwandt.
\pend%
\pstart%
Diese letztere %
stellt ferner das Verhalten der Luft als widerstandsfähigen Mediums in einer Weise dar, welche den Ausführungen über die Übertragung des Schalls in den \textit{Cogitationes novae de sono} bemerkenswert ähnelt (vgl. \textit{LSB} VIII,\;3 N.\;12\textsubscript{1} bis N.\;12\textsubscript{3}, bes. S.\;98.1--99.2, 106.21--109.3 u.\;123.5--126.8).\cite{02527}\cite{01481}\cite{02529}
Diese Ausführungen entstanden allesamt im Zeitraum von der zweiten Hälfte August 1681 bis etwa Ende April 1682 und wurden spätestens in der ersten Hälfte 1685 zum Teil überarbeitet (vgl. die editorische Einleitung zu \textit{LSB} VIII,\;3 N.\;12, S.\;90--93, bes. S.\;90.37--91.27).\cite{02526}
Im gleichen Zeitraum könnte auch -- möglicherweise als Nebenprodukt der Abfassung und Überarbeitung %
der \textit{Cogitationes novae de sono} -- die vorliegende Aufzeichnung N.\;\ref{60349-1} entstanden sein.
Diese könnte aber auch im Zusammenhang mit den ersten Briefen über die physikalische Akustik stehen, die Leibniz bereits im Februar oder März 1681 an Günther Christoph Schelhammer\indexN{\textls{Schelhammer} (Schelhammerus), Günther Christoph 1649\textendash1716} und Edme Mariotte\indexN{\textls{Mariotte}, Edme, Seigneur de Chazeuil ca. 1620\textendash1684} sendete und als Vorlage für die \textit{Cogitationes novae de sono} dienten (siehe erneut die editorische Einleitung zu \textit{LSB} VIII,\;3 N.\;12, S.\;90.4--21).\cite{02526}
\pend%
\pstart%
Die gleich zu Beginn der Aufzeichnung N.\;\ref{60349-1} formulierte These, dass die "Ursache" der Körperkohäsion der \textit{motus corporum ambientium} sei, sofern er Widerstand gegen die Trennung der Körperteile leiste, ist bereits in älteren Texten anzutreffen, etwa in Leibnizens Brief an Honorè Fabri\indexN{\textls{Fabri}, Honoré 1607\textendash1688} wohl von den ersten Monaten 1677 (prop.\;4: \textit{LSB} II,\;1\textsuperscript{2} N.\;133, S.\;452.1--16).\cite{01309}
Wahrscheinlich ist N.\;\ref{60349-1} später entstanden als dieser Brief, %
in dem Leibniz noch die Ruhe als mögliche Ursache der Körperkohäsion zulässt (ebd., %
bes. S.\;452.11\,f.),\cite{01309} während dies in N.\;\ref{60349-1} nicht mehr der Fall ist.
Trotzdem ist nicht auszuschließen, dass N.\;\ref{60349-1} noch vor Leibnizens Auseinandersetzung mit Schelhammer\indexN{\textls{Schelhammer} (Schelhammerus), Günther Christoph 1649\textendash1716} und Mariotte\indexN{\textls{Mariotte}, Edme, Seigneur de Chazeuil ca. 1620\textendash1684} und vor den Texten zur Übertragung des Schalls abgefasst wurde.
Auch eine Entstehung von N.\;\ref{60349-1} nach 1684/1685 ist nicht auszuschließen, wenngleich sie weniger wahrscheinlich ist.
Demgemäß empfiehelt es sich, als Datierung der Aufzeichnung N.\;\ref{60349-1} den gesamten Belegzeitraum des Wasserzeichens zu betrachten.%
\pend%
\selectlanguage{latin}%
\frenchspacing%
\normalsize%
